% ============================================================
\section{Cross-Domain Analysis and Discussion}
\label{sec:analysis}
% ============================================================

% ============================================================
\subsection{Held-Out Scoring Is a Precondition, Not a Refinement}
\label{sec:analysis:heldout}
% ============================================================

Every number in this paper is scored on windows disjoint from the ones used
for early stopping.  That is not a stylistic preference.  Scoring
$\text{B}{-}\text{D}$ on the selection windows instead \textbf{reverses its
sign in 4 of the 31 intervention cells}, and \emph{every one of the four
reverses in the same direction}: from ``the frozen encoder was better'' to
``full fine-tuning was better'' (Appendix Table~\ref{tab:heldout_decomp}).  The largest
is Chronos/ETTh1, which moves from $+6.8$ to $-39.2$, a 46-point swing.  Two
of the four are among the five gate-passing cells of
\S\ref{sec:forecasting:results}, so this is not confined to cells nobody would
have looked at.

The mechanism is visible in the decomposition.  The reversals are driven by
the frozen arm, whose held-out penalty is much larger than its selection-window
penalty ($\Delta$ frozen $+53.9$ against $\Delta$ adapted $+8.0$ on
Chronos/ETTh1), and they concentrate where the two splits are least
exchangeable: the untrained checkpoint's held-out-to-selection error ratio
spans $0.98$ to $5.07$ across these four cells, so on Moirai-Large/ETTh2 the
test split is five times harder than the split used to choose the checkpoint.
A frozen encoder cannot adapt to that shift and a fine-tuned one partly can,
which is exactly the asymmetry that makes same-split scoring flatter freezing.
The prescription is unsurprising; that ignoring it inverts 4 of 31 conclusions,
by up to 46 points and always in the same direction, is not.

% ============================================================
\subsection{Within Moirai, and on Two Further Backbones}
\label{sec:forecasting:nonmoirai}
% ============================================================

\textbf{Drift and outcome come apart inside a single backbone}---within our
Moirai matrix, the strongest form of the finding.  Beyond the sweep of
\S\ref{sec:forecasting:replication}: within Moirai-Base the two \emph{least}
drifted cells are the only two that degrade (ETTh1 at CKA $0.755$ and $0.669$,
$+22.8\%$ and $+39.5\%$) while its two most drifted improve the most within the
arm (ETTh2 at CKA $0.460$ and $0.396$, $-31.4\%$ and $-36.4\%$); and inside
Moirai-Large's LoRA arm the learning rate, not the rank, decides the sign at
near-identical CKA ($0.973$ against $0.989$; Appendix~\ref{app:lora_rank}).

\textbf{Correction to an earlier reading of the non-Moirai gates.} An earlier
version reported ``eight of nine non-Moirai cells gate-pass'' at $h{=}24$;
against the fitted baseline \textbf{none of the ten does}, and the earlier count
was an artefact of the trend denominator (Appendix~\ref{app:gatecorrection}).
What the two extra backbones still supply is the dissociation, which does not
depend on the gate.

\textbf{Near-identical drift, opposite consequences.} On Chronos-T5-Small four
of the five datasets restructure to a degree that would, on the usual reading,
be called catastrophic---CKA ${\in}[0.092, 0.169]$, at or below the
$0.215{\pm}0.054$ an \emph{untrained} encoder of this architecture scores
(Appendix~\ref{app:cka_calibration})---yet outcomes across the five span \textbf{$-41.4\%$ on
ETTm2 to $+20.6\%$ on ETTh1}, a 62-point spread, with ETTh1 the \emph{least}
drifted of the four and the only one that degrades.  TimesFM-2.5 adds four
cells at milder drift (CKA $0.246$--$0.561$), all improving ($-51.8\%$ to
$-3.0\%$).  Pooled across all 31 cells CKA does correlate with the intervention
($\rho{=}{+}0.567$, 95\% CI $[{+}0.265,{+}0.759]$), but that pooled figure is a
backbone artefact by our own rule: within Moirai, the only arm with enough
cells to ask, $\rho{=}{+}0.168$ with CI $[-0.325,{+}0.595]$, which admits a
moderate positive relationship but does not order our cells.  Nor does the
second diagnostic \S\ref{sec:method} defines rescue it: pooled $\ell_2$ weight
drift gives $\rho{=}{+}0.118$ with a CI including zero, on the 25 cells that
retain a drift record, and neither do five further geometric measures on the 59
encoders that retained checkpoints
(Appendices~\ref{app:crosscell},~\ref{app:driftmetrics}).

\textbf{What the frozen control says.} On Chronos/ETTh1 the frozen encoder is
worse still ($+59.8\%$ against full fine-tuning's $+20.6\%$): the cell degrades
whether or not the encoder moves, so encoder adaptation is not the locus of the
damage.

\textbf{ILI: a small-data external case study.}
\label{sec:analysis:ili}
Moirai-Small on ILI (weekly influenza, $h{=}24$, 10 seeds) restructures heavily
(CKA $0.456{\pm}0.008$) and improves ($-50.8{\pm}0.5\%$ held-out, 10/10
negative), but \textbf{gate-fails decisively} at
$R^2_\text{task}(\text{PT}){=}{-}1.625$.  Two further corrections attach to it:
the pre-registered CKA-intermediate prediction ($0.90$--$0.95$) is
\textbf{falsified}, and an earlier report of \emph{zero} frozen-encoder
improvement in all 10 seeds was a script defect, now
$\text{B}{-}\text{D}{=}{-}38.6{\pm}1.3$ (Appendix~\ref{app:ili_gate}).  At
${\approx}450$ windows this is a case study, not a pillar; an IoT
anomaly-detection negative control behaves the same way
(Appendix~\ref{app:cnn_detail}).

% ============================================================
\subsection{No Cell Meets the Degradation Definition}
\label{sec:analysis:degradation}
% ============================================================

Applying the three-clause definition of \S\ref{sec:method} to all 31
intervention cells returns \textbf{zero cells}.  The count comes from
\texttt{scripts/cell\_matrix.py}, not from hand assembly.

Clause~(i) removes 26 of the 31 outright: they have no measured pre-trained
advantage, so a worse post-fine-tuning score is benign replacement rather than
damage.  Of the five that survive it, three fail clause~(ii) in the strongest
possible way---full fine-tuning \emph{improves} MSE by 31--42\% with 3/3 seeds
agreeing.  That leaves two, and both fail the unanimity requirement:
Moirai-Small/ETTh2 at $h{=}192$ has forg.$_\text{B}{=}{+}1.99{\pm}3.32$ with 5
of 10 seeds positive, and at $h{=}96$ the positive mean is a single diverged
seed (\S\ref{sec:forecasting:results}).  The $h{=}192$ cell satisfies
clauses~(i) and~(iii) and is the closest the matrix comes.  Relaxing
unanimity---to a majority, 70\%, 80\%, or a $t$-test, at either gate
threshold---admits at most that one cell (Appendix~\ref{app:defsensitivity}).
Changing the gate's \emph{estimator} instead admits none under three of the four
alternatives and one under a seasonal-naive denominator weaker than ours on that
cell, so loosening the baseline is what manufactures degradation cells
(Appendix~\ref{app:baselines}).

\textbf{What we are and are not saying.}  Not that fine-tuning never damages a
pre-trained capability, but that across 32 screened cells on six standard
benchmarks and three backbones we could not construct one instance that
survives its own definition---fixed before the corrected gate was computed.
These benchmarks are the wrong place to look: on 27 of 32 cells there is no
advantage to lose, and on the five where there is, fine-tuning on
500--1{,}000 samples usually improves on it.
\S\ref{sec:limitations} states what a cell would have to look like to settle
the question.

\textbf{This is a change of conclusion, not of emphasis.}  An earlier version
of this work reported eight degradation cells, all on ETTh1 or Weather at
$h{\in}\{96,192\}$.  All eight fail clause~(i) once the gate is computed
against the linear regression the gate is defined as; none of the eight
involved re-running a model
(Appendix~\ref{app:gatecorrection}).

% ============================================================
\subsection{A Pre-Registered Prospective Test With No Positive Instance}
\label{sec:analysis:prospective}
% ============================================================

Everything above is retrospective: the cells were run, then read.  A
diagnostic that only ever explains cells it has already seen is worth little,
so we tested ours forward.  \textbf{Predictions were written to
\texttt{results/v47\_prospective/preregistration.json} and committed to git at
a named revision before any outcome run for those cells existed}; the scoring
script reads that file and the finished runs and decides nothing.
Two rules were registered against each other:

\begin{itemize}\itemsep1pt \parskip0pt
  \item \textbf{gate rule} --- gate$_\text{val}\!\geq\!0.20 \Rightarrow$ at
    risk of degradation. This is the paper's own criterion.
  \item \textbf{dataset rule} --- dataset ${\in}$ \{ETTh1, Weather\}
    $\Rightarrow$ degradation. A \emph{post hoc} pattern read off the
    original Moirai cells, registered as a competitor so that the gate's
    performance has a reference point rather than being graded against
    nothing.
\end{itemize}

Eight held-out cells were then run (Weather, ETTm2, Electricity7, ETTh1
$\times$ Small/Base/Large, 3 seeds each) and scored by the pre-registered
outcome definition, unchanged (Table~\ref{tab:prospective}).

\begin{table}[t]
\centering
\caption{\textbf{The pre-registered prospective test.} Predictions frozen in
git before any condition-B or -D run for these cells existed. ``Degradation''
applies the paper's own definition unchanged: gate-passing $\wedge$
forg$_\text{B}{>}0$ in every seed $\wedge$ forg$_\text{D}{<}0$ in every seed.}
\label{tab:prospective}
% GENERATED by scripts/score_prospective.py --latex -- do not edit by hand.
\begin{center}
\footnotesize
\setlength{\tabcolsep}{4pt}
\begin{tabular}{@{}lcccccc@{}}
\toprule
Cell (new) & gate$_\text{pre}$ & gate$_\text{val}$ & forg$_\text{B}$ & B$-$D held-out & Degradation? & Predicted \\
\midrule
base Weather ($h{=}96$) & $+0.41$ & $-0.82$ & $-11.7$ (0/3) & $+9.4{\pm}2.0$ & no & gate\,$\times$, data\,$\times$ \\
base Weather ($h{=}192$) & $+0.59$ & $-0.63$ & $-18.1$ (0/3) & $+12.5{\pm}8.7$ & no & gate\,$\times$, data\,$\times$ \\
base ETTm2 ($h{=}96$) & $+0.30$ & $-0.26$ & $-13.1$ (0/3) & $+7.0{\pm}2.4$ & no & gate\,$\times$, data\,$\checkmark$ \\
base ETTm2 ($h{=}192$) & $+0.26$ & $+0.21$ & $-8.4$ (0/3) & $+1.3{\pm}3.4$ & no & gate\,$\times$, data\,$\checkmark$ \\
small Electricity7 ($h{=}96$) & $+0.71$ & $-0.65$ & $-16.0$ (0/3) & $-3.2{\pm}1.5$ & no & gate\,$\times$, data\,$\checkmark$ \\
small Electricity7 ($h{=}192$) & $+0.76$ & $-0.45$ & $-18.6$ (0/3) & $-2.6{\pm}2.1$ & no & gate\,$\times$, data\,$\checkmark$ \\
large ETTh1 ($h{=}96$) & $+0.35$ & $-0.25$ & $+14.1$ (3/3) & $+28.2{\pm}4.9$ & no & gate\,$\times$, data\,$\times$ \\
large Weather ($h{=}96$) & $+0.47$ & $-0.64$ & $+7.9$ (3/3) & $+12.0{\pm}7.5$ & no & gate\,$\times$, data\,$\times$ \\
\bottomrule
\end{tabular}
\end{center}

\end{table}

\textbf{Result: the outcome the gate was registered to predict never occurs.}
None of the eight cells degraded.  The gate flagged all eight, giving
\textbf{TP~0, FP~8, FN~0}: precision $0.00$, and recall undefined because the
denominator is zero.  The competing dataset rule scores no better, for the
same reason.  This is a stronger negative than the one we expected to report,
and it is a different kind: the earlier version of this arm read TP~2, FP~6,
precision $0.25$, recall $1.00$, and we defended the gate as a screen that
never missed a degradation cell.  With the corrected baseline there are no
degradation cells to miss, so the screen has nothing to be right about.

\textbf{The $0.20$ threshold is not what fails.} Re-scoring the same
pre-registered predictor against the same pre-registered outcome at thresholds
from $0.10$ to $0.60$---moving nothing else---gives precision $0.00$ at every
one, with 8, 8, 6, 5, 3 and 2 cells flagged respectively (Appendix
Table~\ref{tab:gate_sensitivity}).  \textbf{The problem is not the cutoff and
it is not even the predictor: it is that the predicted class is empty.}  So the
$0.00$ is arithmetic rather than a measured error rate---with no positives to
find, any rule scores it---and what carries information is the negative
correlation below.

\textbf{The gate is worse than uninformative about the intervention.}
Within Moirai ($n{=}22$ cells from six series, the only arm with enough cells to
ask), the gate's continuous score correlates \emph{negatively} with held-out
$\text{B}{-}\text{D}$: Spearman $\rho{=}{-}0.518$.  Resampling series rather
than cells widens the interval across zero, so we report the direction and not
the interval (Appendix~\ref{app:clustered}).  Higher measured pre-trained value
predicts encoder adaptation helping \emph{more}, not less---the opposite of
what a preservation heuristic needs.  (forg.$_\text{B}$ tracks
$\text{B}{-}\text{D}$ almost perfectly, $\rho{=}{+}0.894$, but only by sharing a
denominator with it---a caution, not a shortcut;
Appendix~\ref{app:predictor}.)

\textbf{What survives.} Not the gate as a predictor of harm, and not as a
screen for when preservation is worth the effort---on these cells it would have
fired eight times out of eight and been wrong eight times.  What survives is
the gate as a \emph{disqualifier}: a cell that fails it cannot supply evidence
about damage to pre-trained features, and 27 of our 32 fail it.  That is a much
weaker instrument than we proposed, and it is why this paper reports a negative
result about the premise rather than a diagnostic for the phenomenon.
